% ===================================================================
%  TABELO N° 7 — La vilajo dum la vintro. La rustika domo dum la
%  printempo.
%  Feuillets 49 a 56 du fac-simile = folios 47 a 54.
%  Correspondance : folio imprime = numero de feuillet - 2.
%
%  Ca tabelo esas en du ceni, kada un kun sua propra numerado
%  d'alineas, ripetanta de 1 : « Unesma ceno » (La vilajo dum la
%  vintro) e « Duesma ceno » (La rustika domo dum la printempo). La
%  cles %%K sequas la numerado quale imprimita. Le feuillet 49 porte
%  anke la ligno d'aparato « DUESMA SERIO » avan la titulo dil tabelo.
% ===================================================================

% --- Feuillet 49 (folio 47, apertajo di la tabelo) -------------------
\begin{VUpage}[49]{}
\VUnotes{143.2mm}{%
(*) Videz pri la repastala detali, la tabeli 6 e 13.%
}
%%K t07-apar-1 apar
\VUsaut{4.0mm}
\VUcentre{13.2pt}{90}{DUESMA SERIO}
\VUsaut{3.0mm}
\VUfilet{16mm}
\VUsaut{5.6mm}
\VUtitre{15.0pt}{60}{TABELO N\textsuperscript{o} 7}
\VUsaut{3.0mm}

%%K t07-tit-1 sub
\VUcentre{12.0pt}{0}{\textit{Unesma ceno.}}
\VUsaut{3.0mm}

%%K t07-tit-2 sub
\VUpk{10.2pt}{40}{La vilajo dum la vintro.}
\VUsaut{4.0mm}

%%K t07-c1-01-1 p
\VUblancAlinea
1. --- Dum la novyaral vakanci, me akom\cc
panis mea patrulo en Novurbo, vilajeto ube\nl
il bezonis traktar kelk aferi komercala.

%%K t07-c1-02-1 p
\VUblancAlinea
2. --- Ni uzis la \VUgras{dilijenco} \textsuperscript{(11)} qua funcionas\nl
inter la urbo e ta burgo; ma pro ke esis tre\nl
kolda, ta voyajo en vehilo poke komfortoza\nl
ne esis tre agreabla.

%%K t07-c1-03-1 p
\VUblancAlinea
3. --- Pro to ni sentis vera plezuro kande ni\nl
vidis la \VUgras{duktisto} haltigar sua veturo sur la\nl
placo, avan l'\VUgras{albergo} \textsuperscript{(2)} di la \textit{Blank Urso}. \VUgras{L'al}\cc
\VUgras{bergisto} \textsuperscript{(4)} esis vartanta ni sur la solio di sua\nl
domo, kalme fumanta sua \VUgras{pipo.}

%%K t07-c1-03-2 p
\VUblancAlinea
Il invitis ni enirar e me ne lasis me prege\cc
sar por instalar me avan la sarmentoza fairo\nl
qua krepitis en la herdo. Lore me tre mokis\nl
l'akra \VUgras{nordal vento} qua grincigas la olda \VUgras{insi}\cc
\VUgras{gno} \textsuperscript{(3)} et forportis per nigra vortici \VUgras{la fumuro} \textsuperscript{(1)}\nl
ekiranta la kameno.

%%K t07-c1-04-1 p
\VUblancAlinea
4. --- Esis la horo dil dejuno. Sen plu longe\nl
vartar, ni bone manjis la simpla ma saporoza\nl
repasto quan on servis a ni (*).
\end{VUpage}

% --- Feuillet 50 (folio 48) ------------------------------------------
\begin{VUpage}[50]{48}
%%K t07-tit-3 sub
\VUpk{10.2pt}{40}{Kelka viziti.}
\VUsaut{3.0mm}

%%K t07-c1-05-1 p
\VUblancAlinea
5. --- Pos la dejuno, me akompanis mea pa\cc
trulo, qua esas asekuristo, che lua klientaro.\nl
Unesme ni vizitis la \VUgras{forjisto} \textsuperscript{(5)}, qua lojas\nl
proxim l'albergo. Il esis okupata per (huf)\cc
ferizar kavalo. Me admiris la robusteso di lua\nl
kompano qua solide mantenis sur lua ledra\nl
avantalo la \VUgras{hufo} dil kavalo, dum ke la for\cc
jisto fixigis la \VUgras{fero} \textsuperscript{(6)} per forta martelagi.

%%K t07-c1-06-1 p
\VUblancAlinea
6. --- Pose ni iris che la \VUgras{shuisto} \textsuperscript{(7)} dil vilajo,\nl
la olda Iakobus. Malgre ke il havas kom insi\cc
gno belega \VUgras{boto,} il saciesis per risuolizar olda\nl
paro de shui truizita.

%%K t07-c1-06-2 p
\VUblancAlinea
La oldo dicis a ni ridante : « En la urbo,\nl
me esis « shuifisto » e me ne havis klienti.\nl
Hike, me esas « shureparisto » e la laboro\nl
esas multa. »

%%K t07-c1-07-1 p
\VUblancAlinea
7. --- Il parolis a ni pri sua vicino la \VUgras{razis}\cc
\VUgras{to} \textsuperscript{(8)}, di qua la klientaro esas ne kontenta. Lua\nl
\VUgras{razili} sempre esas brechizita e lua \VUgras{cizi} nul\cc
tempe tondas bone.

\VUblancAlinea
Ma pro ke il esas la sola razisto dil vilajo,\nl
on esas koaktata komprenende razigar su e\nl
tondigar sua hararo da ilu. Cetere il esas avara\nl
e desafabla viro.

%%K t07-c1-08-1 p
\VUblancAlinea
8. --- Fortunoze, la cetera \VUgras{vicini} ne simi\cc
lesas ilu. Exemple la \VUgras{charifisto} \textsuperscript{(9)} esas beni\cc
gna viro, laborema e komplezema. Esas ple\cc
zuro reparigar veturo da lu. Lua laboruro esas\nl
sempre finisita.

%%K t07-c1-09-1 p
\VUblancAlinea
9. --- Olda Iakobus anke ne plendis pri la\nl
\VUgras{selifisto} \textsuperscript{(10)}, quan il kelkafoye helpas sutar la\nl
\VUgras{harnesi} kande urjas la laboro.
\end{VUpage}

% --- Feuillet 51 (folio 49) ------------------------------------------
\begin{VUpage}[51]{49}
%%K t07-c1-09-2 p
\VUblancAlinea
Mea patrulo questionis lu pri lua familio :\nl
« Omni standas bone. Mea nepotino esas en\nl
la \VUgras{skolo} \textsuperscript{(12)}. Elua \VUgras{docistino} \textsuperscript{(13)} esas tre kon\cc
tenta pri lu, el esas bona \VUgras{lernanto} \textsuperscript{(14)}. El ne\nl
similesas mea mala kerlo filio; il pensas nur\nl
pri ludar. La \VUgras{docisto} \textsuperscript{(15)} ne sucesas lernigar\nl
irgo da ilu. »

%%K t07-tit-4 sub
\VUsaut{3.0mm}
\VUpk{10.2pt}{40}{Vilajala babiladi.}
\VUsaut{3.0mm}

%%K t07-c1-10-1 p
\VUblancAlinea
10. --- La oldo pose rakontis a ni la novaji\nl
dil vilajo. On esas konstruktonta nova skolo,\nl
nam la instalita en la \VUgras{komonal domo} divenis\nl
tro mikra. Cetere to esas facila, nam suficos\nl
komprar parto del gardeno dil \VUgras{muelisto.} To\nl
esos tre profitoza por la kompatinda viro. Pro\nl
la koldeso, la \VUgras{muel-diski} dil \VUgras{mueleyo} \textsuperscript{(16)} ne\nl
plus povas muelar la frumento, pro ke la\nl
\VUgras{roto} \textsuperscript{(17)} esis konjele haltigita da la \VUgras{glacio} \textsuperscript{(25)}\nl
di la \VUgras{rivereto} \textsuperscript{(23)}.

%%K t07-c1-11-1 p
\VUblancAlinea
11. --- « Pri konstrukturo, ka vu vidis nia\nl
nova \VUgras{kirko} \textsuperscript{(18)} ? Ol esas tre pitoreska sub sua\nl
nivura mantelo. La \VUgras{korvi} \textsuperscript{(19)}, qui ne plus ri\cc
konocas sua olda klosheyo, triste perchas sur\nl
la grand arboro dil \VUgras{tombeyo} \textsuperscript{(20)} ».

%%K t07-c1-12-1 p
\VUblancAlinea
12. --- Ta ideo pri la tombeyo rimemorigis\nl
dal shuifisto la personi quin mea patrulo ko\cc
nocabis e qui mortis lastayare. « Quanta\nl
\VUgras{tombi} \textsuperscript{(21)}, mea bon sioro, adjuntesis a la ce\cc
teri, sub la \VUgras{cipreso} \textsuperscript{(22)} ! La morto falchis nia\nl
olda notario. La tota vilajani asistis lua \VUgras{fu}\cc
\VUgras{nero.} »

%%K t07-c1-13-1 p
\VUblancAlinea
13. --- « Yen lua sucedanto qua sketas sur\nl
la rivereto. Il jus venis reparigar da me la\nl
rimeno di sua \VUgras{sketilo} \textsuperscript{(26)} qua esis ruptita. »
\end{VUpage}

% --- Feuillet 52 (folio 50) ------------------------------------------
\begin{VUpage}[52]{50}
%%K t07-c1-14-1 p
\VUblancAlinea
14. --- « Yen anke sur la rivo di la rivereto,\nl
la nova \VUgras{rural gardisto} \textsuperscript{(27)}. Il esas ex-jendarmo\nl
qua terorigas la bubachi dil vilajo. Li fugas\nl
quik kande li videskas lua \VUgras{getri} \textsuperscript{(29)} e lua \VUgras{ke}\cc
\VUgras{pio} \textsuperscript{(28)}. »

%%K t07-c1-15-1 p
\VUblancAlinea
15. --- « Ha ! yen la olda Iosef qua duktas\nl
\VUgras{charetedo de ligna faski} \textsuperscript{(30)} por la panifisto.\nl
--- To esas vera, dicis mea patrulo, me ri\cc
konocas lua jungajo de \VUgras{mulini} \textsuperscript{(31)}. Me ya bezo\cc
nas parolar ad ilu, me profitos l'okaziono. Til\nl
rivido, patro Iakobus. »

%%K t07-c1-16-1 p
\VUblancAlinea
16. --- Jus kande ni ekiris, la pueri dil\nl
vilajo livis la skoleyo e kurante precipitis su\nl
ad la \VUgras{gliteyo} \textsuperscript{(33)} kun risko falar e ruptar a su\nl
membro en la \VUgras{falo} \textsuperscript{(34)}. Un de li prenis sua\nl
\VUgras{glit-veturo} \textsuperscript{(32)} che la charifisto, sidigis en olu\nl
un de sua kamaradi, e departis kurante.

%%K t07-c1-17-1 p
\VUblancAlinea
17. --- Altri precipitis su ad \VUgras{niv-amaso} \textsuperscript{(37)}\nl
proxim la \VUgras{ponto} \textsuperscript{(40)} e bombardeskis la \VUgras{nivura}\nl
\VUgras{statuo} \textsuperscript{(39)} quan li erektis dum la hiera dio ed\nl
a qua li metis chapelo, pipo e bandoliera\nl
skolsako.

%%K t07-c1-18-1 p
\VUblancAlinea
18. --- Li poke sucias la frostiganta vento e\nl
la koldeso. La plumulto mem ne havas \VUgras{kra}\cc
\VUgras{vatego} \textsuperscript{(35)} e, por rivarmigar su pos la kom\cc
bato per \VUgras{niv-buli} \textsuperscript{(38)}, suficas a li manuedo de\nl
tre varma \VUgras{kastani} \textsuperscript{(42)} komprita de la olda\nl
\VUgras{vendistino} Jakobina, qua tremas pro koldeso\nl
sub sua \VUgras{shalacho} \textsuperscript{(43)} tro dina.
\end{VUpage}

% --- Feuillet 53 (folio 51) ------------------------------------------
\begin{VUpage}[53]{51}
%%K t07-tit-5 sub
\VUsaut{3.0mm}
\VUcentre{12.0pt}{0}{\textit{Duesma ceno.}}
\VUsaut{3.0mm}

%%K t07-tit-6 sub
\VUpk{10.2pt}{40}{La rustika domo dum la printempo.}
\VUsaut{4.0mm}

%%K t07-c2-01-1 p
\VUblancAlinea
1. --- Malgre omna plezuri dil vintro, ni kun\nl
profunda joyo salutas la riveno dil \VUgras{printempo.}

%%K t07-c2-01-2 p
\VUblancAlinea
Quante gayiganta tabelon aspektas farmala\nl
korto, vespere, en la mayo-monato !

%%K t07-c2-02-1 p
\VUblancAlinea
2. --- Ye la horizonto la \VUgras{suno} \textsuperscript{(1)} lente kushas\nl
su, inundante la \VUgras{ruro} \textsuperscript{(3)} per sua purpurea \VUgras{ra}\cc
\VUgras{dii} \textsuperscript{(2)}. La diligenta \VUgras{plugisto} \textsuperscript{(6)} hastas plugar\nl
sua \VUgras{agro} \textsuperscript{(4)}.

%%K t07-c2-02-2 p
\VUblancAlinea
Per sua \VUgras{plugilo} \textsuperscript{(5)} jungita a vigoroza \VUgras{ka}\cc
\VUgras{valo} \textsuperscript{(7)}, il trasas la \VUgras{sulko} \textsuperscript{(8)} en qua la \VUgras{se}\cc
\VUgras{mero} \textsuperscript{(9)} per plena manuedo jetos la \VUgras{semino}\nl
qua produktos la orea spiki.

%%K t07-c2-03-1 p
\VUblancAlinea
3. --- La \VUgras{falchisti} \textsuperscript{(11)} provizita per sua fal\cc
chili falchis la herbo dil prato, la fanigisti\nl
sikigis la \VUgras{feno} \textsuperscript{(10)} turnante lu per \VUgras{forki} \textsuperscript{(12)} e\nl
rastili, e yen ke li retrovenas.

%%K t07-c2-03-2 p
\VUblancAlinea
Uni marchas avan la pezoza veturo tirata da\nl
bovi; li kunportas sua utensilo sur la shultro.\nl
Altri, plu indolenta, komfortoze su instalis sur\nl
l'odoranta feno.

%%K t07-c2-04-1 p
\VUblancAlinea
4. --- Pasante li facas amikal saluto al \VUgras{gar}\cc
\VUgras{denisto} \textsuperscript{(16)} okupata en la \VUgras{frukto-gardeno} \textsuperscript{(13)}\nl
per taliar la \VUgras{frukt-arbori} e greftar la \VUgras{pirieri} \textsuperscript{(14)}.\nl
La \VUgras{prunieri} \textsuperscript{(15)} florifeskas, e, en la \VUgras{legum-gar}\cc
\VUgras{deno} \textsuperscript{(17)} bone sterkizita, la legumi, kauli e\nl
saladi ja esas bonastanda.

%%K t07-c2-04-2 p
\VUblancAlinea
Sur la mureto, bela \VUgras{pavono} \textsuperscript{(18)} developas\nl
sua kaudo, por admirigar olua plumi belega.
\end{VUpage}

% --- Feuillet 54 (folio 52) ------------------------------------------
\begin{VUpage}[54]{52}
%%K t07-tit-7 sub
\VUpk{10.2pt}{40}{Farmo-korto.}
\VUsaut{3.0mm}

%%K t07-c2-05-1 p
\VUblancAlinea
5. --- La pordi di la \VUgras{kavaleyo} \textsuperscript{(19)} esas aper\cc
tita ed on vidas la rastelero e la \VUgras{litiero} \textsuperscript{(23)} di\nl
la \VUgras{kavalino} \textsuperscript{(21)} quan la \VUgras{kaval-servisto} \textsuperscript{(20)} jus\nl
desjungis. La \VUgras{kavalyuno} \textsuperscript{(22)} tante komika kun\nl
sua longa gambi, sequas sua patrino gambo\cc
lante.

%%K t07-c2-06-1 p
\VUblancAlinea
6. --- Proxim la \VUgras{puteo} \textsuperscript{(29)} la \VUgras{bovini} \textsuperscript{(25)} iras\nl
drinkar en la ligna \VUgras{trogo} \textsuperscript{(26)} qua uzesas por li\nl
kom drinkeyo. La \VUgras{bovyuno} \textsuperscript{(24)} profitas l'oka\cc
ziono por mamsugar, e la \VUgras{bovulo} \textsuperscript{(27)} flarante\nl
la bono doro dil forejo, joyoze bramas e fiere\nl
levas la kapo kun \VUgras{korni} \textsuperscript{(28)} potenta.

%%K t07-c2-07-1 p
\VUblancAlinea
7. --- Omni laboras. La \VUgras{servistino} \textsuperscript{(32)} ma\cc
nue tenas la \VUgras{kordo} \textsuperscript{(31)} dil puteo. El cherpas\nl
aquo per \VUgras{sitelo} \textsuperscript{(30)} por drinkigar la brutaro.\nl
Proxim la \VUgras{garbeyo} \textsuperscript{(33)}, viro rastas la paliaji\nl
ed avan la pordo di la \VUgras{hemo} \textsuperscript{(34)} olda \VUgras{filifis}\cc
\VUgras{tino} \textsuperscript{(35)} per sua filif-roto e sua \VUgras{ruko,} filifas la\nl
\VUgras{lino} o la \VUgras{kanabo} quale en la tempo olima.

%%K t07-tit-8 sub
\VUsaut{3.0mm}
\VUpk{10.2pt}{40}{La farmisto.}
\VUsaut{3.0mm}

%%K t07-c2-08-1 p
\VUblancAlinea
8. --- La \VUgras{farmisto} \textsuperscript{(36)} direktas ta omna labori\nl
e donas instrucioni a sua servisti. Il reko\cc
mendas subtektigar la \VUgras{herso} \textsuperscript{(40)} e la \VUgras{piocho} \textsuperscript{(39)}\nl
quin on obliviis proxim la \VUgras{kabano} \textsuperscript{(38)} di la\nl
\VUgras{hundo} \textsuperscript{(37)}.

%%K t07-c2-09-1 p
\VUblancAlinea
9. --- Lua klami pavorigas \VUgras{kolombo} \textsuperscript{(41)}, qua\nl
per tota ali flugas aden la \VUgras{kolombeyo} \textsuperscript{(42)}, e la\nl
\VUgras{hanini} \textsuperscript{(43)} qui hastas rienirar la \VUgras{haneyo} \textsuperscript{(44)}.

%%K t07-c2-10-1 p
\VUblancAlinea
10. --- Avan la \VUgras{mutoneyo} \textsuperscript{(48)}, yen la olda\nl
\VUgras{pastoro} \textsuperscript{(45)} qua retroduktas sua (muton)-\VUgras{tru}\cc
\parplein
\end{VUpage}

% --- Feuillet 55 (folio 53) ------------------------------------------
\begin{VUpage}[55]{53}
%%K t07-c2-10-1 p suite
\VUcontinue
\VUgras{po} \textsuperscript{(49)}. Tenante sua (pastor)-\VUgras{bastono} \textsuperscript{(46)}, qua\nl
nultempe mankas a lu, same kam lua \VUgras{bluzo} \textsuperscript{(47)}\nl
senkolorigita, il duktas a la stablo \VUgras{muto}\cc
\VUgras{nuli} \textsuperscript{(50)}, \VUgras{mutonini} \textsuperscript{(53)}, \VUgras{mutonyuni} \textsuperscript{(54)}, \VUgras{kapri}\cc
\VUgras{ni} \textsuperscript{(51)} e \VUgras{kaproyuni} \textsuperscript{(52)}. Lua fidela \VUgras{hundo} \textsuperscript{(55)}\nl
Argus laste venas ed instigas per sua aboyi la\nl
diligenteso dil tardeskanti.

%%K t07-c2-11-1 p
\VUblancAlinea
11. --- Ta omna bruisado e movado ne tru\cc
blas la quieteso dil bona \VUgras{laktifanta bovino} \textsuperscript{(56)}\nl
qua tinkligas sua \VUgras{klosheto} \textsuperscript{(57)}, dum ke on mel\cc
kas elu avan la pordo dil \VUgras{stablo} \textsuperscript{(58)}.

%%K t07-c2-12-1 p
\VUblancAlinea
12. El semblas komprenar, ke el devas do\cc
nar sua lakto, por ke la \VUgras{farmistino} \textsuperscript{(60)} povez\nl
prontigar \VUgras{butro} en la \VUgras{butrifilo} \textsuperscript{(61)} o la bonega\nl
\VUgras{fromaji} \textsuperscript{(64)}, qui gutifas de la tabulo pozita sur\nl
la \VUgras{kuvego} \textsuperscript{(63)}.

%%K t07-c2-13-1 p
\VUblancAlinea
13. --- La tota \VUgras{lakteyo} \textsuperscript{(59)} mantenesas tre\nl
nete da la \VUgras{farm-servisto} \textsuperscript{(65)} qua jus lavis la\nl
\VUgras{lakto-poti} \textsuperscript{(62)} en la granda sitelo quan il portas\nl
per la manuo.

%%K t07-tit-9 sub
\VUsaut{3.0mm}
\VUpk{10.2pt}{40}{La pultreyo.}
\VUsaut{3.0mm}

%%K t07-c2-14-1 p
\VUblancAlinea
14. --- Kun sua dika ligno-shui, il marchas\nl
kelke plumpe, e tale fugigas la \VUgras{dindulo} \textsuperscript{(68)}, la\nl
\VUgras{anadini} \textsuperscript{(66)}, la \VUgras{anaduli} \textsuperscript{(67)} e la \VUgras{gansi} \textsuperscript{(69)} qui\nl
larje apertas \VUgras{sua beko} \textsuperscript{(70)} e desquiete kriegas.

%%K t07-c2-14-2 p
\VUblancAlinea
La \VUgras{hanulo} \textsuperscript{(71)}, plu kombatema, erektas sua\nl
reda \VUgras{kresto} \textsuperscript{(72)}, sukusas la plumi di sua\nl
\VUgras{kaudo} \textsuperscript{(73)} ed audigas « kokoriko » sonorega.

%%K t07-c2-15-1 p
\VUblancAlinea
15. --- \VUgras{Hanin-matro} \textsuperscript{(74)} marodas sur la\nl
\VUgras{sterko} \textsuperscript{(81)} kun sua \VUgras{(han-)yuneti} \textsuperscript{(75)}. La \VUgras{pork}\cc
\VUgras{yuno} \textsuperscript{(78)} fugas ad sua matro, l'enorma \VUgras{por}\cc
\VUgras{kino} \textsuperscript{(77)} quan ni vidas proxim la porkeyo.
\end{VUpage}

% --- Feuillet 56 (folio 54) ------------------------------------------
\begin{VUpage}[56]{54}
%%K t07-c2-16-1 p
\VUblancAlinea
16. --- En sua faki, la \VUgras{kunikli} \textsuperscript{(79)} manjas\nl
avide la \VUgras{herbo} \textsuperscript{(80)} delikata quan adportas a li\nl
la filiino dil farmistulo. Ioanina tre amas sua\nl
kunikli e, singladie, querinta la fresha \VUgras{ovi} \textsuperscript{(76)}\nl
en la haneyo, el venas vizitar sua amiketi.

\VUsaut{6.0mm}
\VUfilet{22mm}
\end{VUpage}
